% ============================================================
% Compile with XeLaTeX (system fonts are used).
%   xelatex proposal.tex ; bibtex proposal ; xelatex ; xelatex
% ============================================================
% !TeX program = xelatex
% !TeX root = proposal.tex
%
% UFS PhD research proposal. Section structure and page budgets follow the
% BITM / UFS Department of Computer Science & Informatics proposal template
% (v4). Section 9 is restructured around the actual work phases, because this
% proposal develops and then tests statistical guarantees rather than running
% an applied data-mining cycle; the deviation is flagged in that section.
%
% Build the cover first (one page):
%   cd ../cover && xelatex cover.tex

\PassOptionsToPackage{final}{graphicx}
\documentclass[12pt,final]{article}

\usepackage{msc_body_style}
\usepackage{changepage}

\begin{document}

	% =============================================================
	% Cover / title page (built separately, included as PDF)
	% =============================================================
	\includepdf[pages=1]{../cover/cover.pdf}

	\clearpage
	\pagenumbering{roman}

	\renewcommand{\contentsname}{Table of Contents}
	\tableofcontents
	\newpage

	\renewcommand{\listfigurename}{List of Figures}
	\listoffigures

	\renewcommand{\listtablename}{List of Tables}
	\listoftables
	\newpage

	\section*{List of Abbreviations}
	\addcontentsline{toc}{section}{List of Abbreviations}

	\begin{tabular}{@{}ll@{}}
		CRPS  & Continuous ranked probability score \\
		DoE   & Design of experiments \\
		DNN   & Deep neural network \\
		GMAB  & Guaranteed minimum accumulation benefit \\
		GMDB  & Guaranteed minimum death benefit \\
		GMWB  & Guaranteed minimum withdrawal benefit \\
		GMxB  & Guaranteed minimum benefit (generic rider) \\
		GRU   & Gated recurrent unit \\
		HMD   & Human Mortality Database \\
		LSMC  & Least-squares Monte Carlo \\
		MC    & Monte Carlo \\
		NLL   & Negative log-likelihood \\
		PV    & Present value \\
		UQ    & Uncertainty quantification \\
		VA    & Variable annuity \\
	\end{tabular}

	\newpage

	\section*{Glossary}
	\addcontentsline{toc}{section}{Glossary}

	\begin{description}
		\item[Aleatoric uncertainty] The irreducible component of predictive error, here the Monte Carlo noise of the simulated labels and the genuine financial and mortality randomness.
		\item[Emulator (surrogate)] A fast approximating function, learned from expensively computed examples, that replaces a costly valuation routine.
		\item[Epistemic uncertainty] The reducible component of predictive error attributable to the emulator itself: its approximation and estimation error.
		\item[Marginal coverage] A coverage probability averaged over calibration samples, as distinct from coverage conditional on the realised calibration sample.
		\item[Nested Monte Carlo] Valuation by an outer layer of scenarios, each carrying an inner layer of risk-neutral paths; the industry reference for guarantee valuation.
		\item[Predictive interval] An interval reported alongside a point prediction that is intended to contain the true target with a stated probability.
		\item[Split-conformal prediction] A distribution-free procedure that calibrates prediction intervals on held-out residuals to achieve finite-sample coverage under exchangeability.
		\item[Stochastic mortality] A model in which the force of mortality is itself a stochastic process rather than a fixed life-table value.
	\end{description}

	\newpage
	\pagenumbering{arabic}

	% ============================================================
	\section{Introduction} % Approx. 1.5 pages

	A variable annuity (VA) combines investment in a reference fund with a set of
	insurance guarantees. The policyholder pays a premium that is invested on
	their behalf, and in return the insurer promises a floor on the benefit paid
	at maturity, on death before maturity, or on a stream of withdrawals, namely
	the guaranteed minimum accumulation, death and withdrawal benefits (GMAB,
	GMDB, GMWB). These guarantees are options written on the policyholder's fund,
	funded by a fee deducted continuously from the account value, and they
	transfer to the insurer a liability that is simultaneously exposed to equity,
	interest-rate and mortality risk. Because the guarantee value has no closed
	form once realistic dynamics are assumed, it must be computed as a discounted
	expectation under an appropriate measure, and the industry reference for that
	computation is nested Monte Carlo simulation \citep{BauerKlingRuss2008}.

	Nested simulation is expensive by construction. An outer layer of scenarios
	describes how the market and the account might evolve, and at each outer node
	an inner layer of risk-neutral paths values the remaining guarantee. The cost
	is the product of the two layers, and it grows sharply when the valuation
	state is high-dimensional. Making the mortality basis itself stochastic, so
	that the force of mortality is a diffusion or a longer-memory process rather
	than a fixed life-table number, enlarges the state further and couples the
	actuarial and financial sides of the problem. For an insurer that must revalue
	a large, heterogeneous book every day for pricing, that stresses it under many
	scenarios for solvency capital, and that re-hedges it as the market moves,
	direct nested simulation is simply too slow.

	Emulation, also called metamodelling or surrogate modelling, is the
	established remedy. Rather than re-solving the valuation from scratch for every
	contract and every state, one learns the valuation function once, from a
	designed set of expensively computed examples, and then evaluates that learned
	function cheaply wherever it is needed. Neural networks are a natural choice of
	emulator because they approximate high-dimensional nonlinear maps without
	hand-crafted basis functions, and because one trained network can value a
	whole portfolio in a single forward pass. The work of \citet{HejaziJackson2016}
	and the metamodelling programme of \citet{GanLin2015, GanLin2018} established
	that this approach can value large VA portfolios to actuarial accuracy at a
	fraction of the cost of nested simulation.

	What those emulators return, however, is a single number. For pricing that may
	suffice; for risk and capital it does not. A valuation that drives a reserve or
	a hedge needs a statement of how far it can be trusted, that is, a quantified
	and calibrated uncertainty that separates the error introduced by the surrogate
	itself from the irreducible randomness of the financial and mortality world.
	Supplying that calibrated uncertainty, under a stochastic-mortality model, is
	the purpose of this research. Concretely, the project trains a deep neural
	network to map a contract, market and mortality state vector directly to the
	present value of the guarantee liability, and then equips every prediction with
	an interval that is guaranteed, in a precise finite-sample sense, to contain the
	true value with a stated probability.

	The project also builds on the candidate's MSc, which developed deep-learning
	methods for hedging path-dependent claims in incomplete markets, combining
	feedforward and recurrent architectures with path-signature features. The
	present work reuses that craft, namely network design, simulation-based
	training and feature construction, but redirects it from hedging toward
	emulation and uncertainty quantification for insurance guarantees that carry
	mortality risk.

	% ============================================================
	\section{Background} % Approx. 3 pages

	\textbf{Variable-annuity guarantees and nested valuation.}
	\citet{BauerKlingRuss2008} set out a unifying framework in which the various
	guaranteed minimum benefits are valued within a single nested-simulation
	architecture. A typical contract invests the premium in a fund of value
	$F_t$, deducts a continuous fee at rate $c$, and promises a guarantee level
	$G$. If the policyholder survives to maturity $T$ the payoff is $\max(G,F_T)$;
	if death occurs at time $\tau<T$ the death benefit $\max(G,F_\tau)$ is paid;
	and if the contract is surrendered the policyholder receives the account value
	net of a surrender charge. The guarantee liability is the value of the shortfall
	options embedded in these promises, and it is this liability, not the account
	itself, that the insurer must reserve against. This framework fixes the
	contract conventions (fee deduction, guarantee bases, surrender-charge
	schedules and withdrawal rules) that define the inputs of any emulator and
	underlies the ground-truth labels used in this project.

	\textbf{Why nested simulation is costly.} With $N_{\text{out}}$ outer scenarios
	and $N_{\text{in}}$ inner paths each, a single portfolio revaluation costs on
	the order of $N_{\text{out}}\,N_{\text{in}}$ path simulations per contract.
	Capital calculations that require the full loss distribution, and daily
	revaluation of a book of hundreds of thousands of policies, multiply this cost
	again. Introducing stochastic mortality adds state variables and, for
	longer-memory specifications, destroys the Markov property, so that standard
	variance-reduction and dynamic-programming shortcuts no longer apply directly.

	\textbf{Emulation and metamodelling of VA portfolios.}
	\citet{HejaziJackson2016} introduced a neural-network approach to the efficient
	valuation of large VA portfolios, learning the valuation map from a
	representative set of contracts and extending it across the portfolio, and
	\citet{GanLin2015, GanLin2018} developed a complementary metamodelling
	programme, including functional-data and two-level designs for both values and
	Greeks. This literature establishes the feasibility and accuracy of emulation,
	but it returns point estimates and typically assumes a static mortality basis
	and no quantified uncertainty.

	\textbf{Deep learning solvers for VA valuation.} Two 2026 contributions are the
	closest prior art. \citet{LiLyu2026} price a VA with an early-surrender option
	under a rough-Heston equity model and a Volterra (long-range-dependent)
	stochastic-mortality model using a deep signature least-squares Monte Carlo
	scheme, and supply a convergence proof that decomposes the pricing error into
	Monte Carlo, time-discretisation, signature-truncation and
	neural-network-approximation parts. \citet{LangreneEtAl2026} extend deep LSMC to
	the guaranteed minimum withdrawal benefit under optimal dynamic withdrawal,
	compare polynomial with neural-network regression, and show how to obtain
	confidence intervals for the contract value. Both are per-contract solvers that
	regress continuation values backward through time; the uncertainty they report
	is the Monte Carlo estimator error of a single price, not a predictive interval
	for an emulated valuation across the state space.

	\textbf{Stochastic mortality.} The Lee-Carter model \citep{LeeCarter1992}
	remains the reference for mortality forecasting, while affine
	stochastic-intensity models \citep{Biffis2005} provide the continuous-time,
	tractable dynamics most convenient for coupling with financial simulation. In
	an affine specification the force of mortality $\lambda_t$ follows a diffusion
	whose survival factor over $[t,s]$ has the exponential-affine form
	\begin{equation}
		{}_{s-t}p_{x+t}
		= \mathbb{E}\!\left[\exp\!\Big(-\!\int_t^s \lambda_u\,du\Big)\,\Big|\,\mathcal{F}_t\right]
		= \exp\!\big(\alpha(s-t) + \beta(s-t)\,\lambda_t\big),
	\end{equation}
	with $\alpha,\beta$ solving the associated Riccati system. Longer-memory
	Volterra specifications, as calibrated to the HMD by \citet{LiLyu2026}, offer a
	richer data-generating process that this project can adopt for its labels.

	\textbf{Uncertainty quantification in deep learning.} Deep ensembles
	\citep{Lakshminarayanan2017} give a simple and strong estimate of epistemic
	uncertainty by training several networks from different initialisations;
	Monte-Carlo dropout interprets dropout at test time as approximate Bayesian
	inference \citep{GalGhahramani2016}; and the separation of predictive
	uncertainty into aleatoric and epistemic parts is formalised by
	\citet{KendallGal2017}. These supply an uncertainty representation but, on their
	own, carry no finite-sample guarantee that the reported interval actually covers
	the truth.

	\textbf{Conformal prediction.} The conformal framework \citep{Vovk2005} and its
	split-conformal regression form \citep{Lei2018} produce prediction intervals
	with distribution-free, finite-sample coverage under exchangeability. Given a
	point predictor and a held-out calibration set, the empirical quantile of a
	nonconformity score yields an interval whose marginal coverage is at least the
	nominal level, for any data distribution and any sample size. The synthesis
	pursued here, a neural VA emulator whose predictive intervals are both
	uncertainty-decomposed and conformally certified under stochastic mortality,
	sits at the intersection of these strands and is, to the candidate's knowledge,
	unaddressed.

	% ============================================================
	\section{Research Gaps} % Approx. 0.5 page

	The gap must be stated carefully against the 2026 deep-LSMC work. Those methods
	are solvers: they price a single contract by backward regression and quantify
	only the Monte Carlo estimator error of that one price. The cross-sectional
	emulators of \citet{HejaziJackson2016} and \citet{GanLin2015}, in turn, return
	point values under static mortality and carry no uncertainty at all. No existing
	method delivers a VA-guarantee emulator that (i) ingests a stochastic-mortality
	state as a first-class input, (ii) returns a distribution-free, finite-sample
	coverage-guaranteed predictive interval for the guarantee value, and (iii)
	decomposes that interval into epistemic and aleatoric components. Rather than
	competing with the LSMC solvers, this work uses them, and in particular the
	calibrated stochastic-mortality dynamics of \citet{LiLyu2026}, as the nested
	Monte Carlo ground truth that generates the emulator's labels and as the
	accuracy benchmark.

	% ============================================================
	\section{Research Problem and Statement} % Approx. 1 page

	A fast emulator that returns only a point value cannot, on its own, support a
	reserve or a solvency figure, because it is silent about its own error. That
	error has two very different sources. The first is the surrogate's own
	approximation and estimation error, which is epistemic and reducible with more
	training data and capacity. The second is the irreducible noise of the
	simulated labels and of the financial and mortality world, which is aleatoric.
	Reporting a single number conflates the two and offers no guarantee that the
	true value lies anywhere near the prediction, which is precisely what a capital
	calculation needs to know. The problem this proposal addresses is therefore to
	construct a VA-guarantee emulator under stochastic mortality whose output is an
	interval with a certified probability of containing the true value, together
	with an honest decomposition of where the uncertainty comes from and how it
	shrinks as computational effort increases.

	% ============================================================
	\section{Thesis Statement (PhD only)} % Approx. 1-2 sentences

	A deep neural network can emulate the present value of a variable-annuity
	guarantee liability under stochastic mortality to an accuracy comparable with
	the Monte Carlo error of the labels it is trained on, and its predictions can be
	equipped with distribution-free, finite-sample coverage-guaranteed intervals
	whose epistemic and aleatoric components are separately identified.

	% ============================================================
	\section{Research Aim} % Approx. 0.5 page

	To develop and evaluate a calibrated, uncertainty-aware deep-learning emulator
	for variable-annuity guarantee valuation under stochastic mortality, and to
	establish the statistical guarantees that make its predictive intervals
	trustworthy for pricing, hedging and capital.

	% ============================================================
	\section{Research Questions and Objectives} % Approx. 1.5 pages

	\textbf{RQ1.} Can a deep neural network emulate the present value of a VA
	guarantee liability, as a function of the full policy, financial-state and
	mortality state vector, to an accuracy comparable with the Monte Carlo error of
	nested-simulation targets, across the relevant ranges of moneyness, age and
	maturity?

	\textbf{RQ2.} How should the emulator's predictive uncertainty be represented
	and decomposed into an epistemic component, attributable to the surrogate, and
	an aleatoric component, attributable to the labels and to financial and
	mortality randomness?

	\textbf{RQ3.} Can the predictive intervals be endowed with a finite-sample
	coverage guarantee through split-conformal prediction, and under what conditions
	does the required exchangeability hold given the way the labels are generated,
	including the effect of Monte Carlo label noise on coverage of the true value?

	\textbf{RQ4.} What is the practical effect of the quantified uncertainty on
	reserves and solvency capital, demonstrated on realistic, data-calibrated
	contracts, and how does interval width trade off against the computational
	budget spent on labels?

	\medskip
	\textbf{Objective 1.} Build the nested Monte Carlo label generator and train the
	emulator, documenting its accuracy against held-out nested simulation across the
	state space.

	\textbf{Objective 2.} Develop and compare uncertainty-quantification schemes,
	namely deep ensembles, Bayesian last-layer approximations and a learned-variance
	head, and separate epistemic from aleatoric error.

	\textbf{Objective 3.} Prove and empirically verify coverage of the conformal
	predictive intervals, accounting for the Monte Carlo noise in the labels.

	\textbf{Objective 4.} Conduct a capital-impact study on HMD calibrations, with
	an optional South African case study, and quantify the label-budget to
	interval-width trade-off.

	% ============================================================
	\section{Hypotheses (Optional)} % Approx. 0.5 page

	These are stated as falsifiable conjectures, not results. \textbf{H1:} a
	feedforward or recurrent emulator attains a relative valuation error against
	nested Monte Carlo comparable to the inner-simulation standard error, uniformly
	over the sampled state space. \textbf{H2:} because the training states are
	sampled independently, exchangeability holds by construction and split-conformal
	intervals achieve their nominal marginal coverage of the labels; coverage of the
	true value follows once label noise is accounted for, either by inflating the
	inner-path count or by subtracting the known label variance.

	% ============================================================
	\section{Research Design and Methodology} % Approx 5-6 pages

	\subsection{Introduction}
	The source template introduces the research design through Saunders' research
	onion and a CRISP-DM cycle, which suit applied information-systems and
	data-mining work. This proposal instead develops statistical guarantees and then
	tests them, so the design is stated directly and the techniques are organised
	around the actual work phases. The deviation is deliberate and is flagged here.

	\subsection{Research Philosophy}
	The work is positioned within a quantitative, model-based tradition. Results are
	derived from stated assumptions and then checked against controlled simulation
	and data, with risk-neutral pricing assumptions kept separate from historical
	estimation of the mortality and market models.

	\subsection{Approach to Theory Development}
	Deductive. Coverage and decomposition results are derived from explicit
	assumptions and then verified numerically. Simulation supports a proof; it does
	not replace one.

	\subsection{Methodological Choice}
	Quantitative throughout: mathematical analysis of the uncertainty guarantees,
	controlled simulation against nested Monte Carlo ground truth, and an empirical
	application calibrated to published mortality data.

	\subsection{Research Strategy}
	Construction of the emulator and its uncertainty wrapper, followed by a coverage
	analysis with known data-generating laws, followed by an application limited to
	what the mortality and contract data support.

	\subsection{Time Horizon}
	Longitudinal in the empirical component: training, calibration and evaluation
	periods are separated in time so that reported coverage reflects genuine
	out-of-sample behaviour rather than in-sample fit.

	\subsection{Problem Formulation}
	Let the financial market carry an equity index $S_t$ and a short rate $r_t$
	under a risk-neutral measure $\mathbb{Q}$, and let the policyholder's death time
	$\tau$ have a stochastic intensity $\lambda_t$ under the physical measure
	$\mathbb{P}$. As is standard in the actuarial literature, and following
	\citet{Biffis2005} and \citet{LiLyu2026}, financial and mortality risks are
	taken to be independent given their drivers, so valuation is carried out under
	the product measure, written $\mathbb{Q}$ for brevity. The account value evolves
	as $F_t = S_t e^{-ct}$ under a continuous fee $c$, and the guarantee generates a
	discounted cashflow stream $\Pi$ determined by the maturity, death and surrender
	benefits described in Section 2.

	Collect in a single state vector $x\in\mathcal{X}\subset\mathbb{R}^d$ the three
	feature blocks that determine the liability:
	\emph{policy characteristics} (age; sex, where it enters the mortality basis;
	contract duration; initial premium; current account value; guarantee base; time
	to maturity; surrender-charge schedule; withdrawal rate; guarantee type),
	\emph{financial-state variables} (interest rate; equity level or moneyness;
	equity volatility; stochastic-volatility state; correlation parameters; fee
	rate) and \emph{mortality variables} (current intensity; trend; volatility;
	cohort effect; financial-mortality correlation; mortality-model parameters). The
	quantity to be emulated is the present value of the guarantee liability
	conditional on the state,
	\begin{equation}
		Y(x) \;=\; \mathbb{E}\!\left[\, \Pi \;\middle|\; \text{state}=x \,\right],
		\label{eq:target}
	\end{equation}
	a deterministic function of $x$ which has no closed form and must be approached
	numerically.

	\subsection{Techniques and Procedures}
	The phases below replace the template's CRISP-DM cycle.

	\subsubsection{Label Generation by Nested Monte Carlo}
	For a state $x_i$ drawn from a design over $\mathcal{X}$, an inner simulation of
	$M$ risk-neutral paths under the chosen stochastic-mortality model produces a
	label
	\begin{equation}
		\widehat{Y}_i \;=\; \frac{1}{M}\sum_{j=1}^{M} \Pi_{ij},
		\qquad
		\mathbb{E}\!\left[\widehat{Y}_i \mid x_i\right] = Y(x_i),
		\qquad
		\operatorname{Var}\!\left[\widehat{Y}_i \mid x_i\right] = \frac{\sigma^2(x_i)}{M}.
		\label{eq:label}
	\end{equation}
	The label is thus an unbiased but noisy estimate of the true target
	\eqref{eq:target}, with a known, controllable noise level. The distinction
	between the deterministic target $Y(x)$ and the noisy label $\widehat{Y}$ is
	central to the coverage analysis below. Path-dependent contracts are summarised
	by signature features of the simulated trajectories, following \citet{LiLyu2026}.

	\subsubsection{Emulator Architecture and Training}
	The emulator $f_\theta:\mathcal{X}\to\mathbb{R}$ is trained on
	$\{(x_i,\widehat{Y}_i)\}_{i=1}^{n}$ by minimising a regression loss
	\begin{equation}
		\mathcal{L}(\theta) \;=\; \frac{1}{n}\sum_{i=1}^{n} w(x_i)\,
		\big(f_\theta(x_i) - \widehat{Y}_i\big)^2,
	\end{equation}
	with optional weights $w$ to balance regions of the state space. Feedforward
	networks form the baseline, with gated-recurrent (GRU) variants where the
	contract state is naturally sequential. Regularisation and hyperparameters are
	selected by Bayesian optimisation, and accuracy is reported against held-out
	nested-simulation values across moneyness, age, maturity and mortality regime.

	\subsubsection{Uncertainty Quantification}
	Epistemic uncertainty is estimated by a deep ensemble of $K$ networks
	\citep{Lakshminarayanan2017}, giving a predictive mean $\bar f(x)$ and an
	epistemic variance $s^2_{\mathrm{epi}}(x)$ equal to the ensemble variance;
	Bayesian last-layer and Monte-Carlo-dropout approximations \citep{GalGhahramani2016}
	are compared. Aleatoric uncertainty is modelled by a variance head
	$\sigma^2_\phi(x)$ trained under the heteroscedastic Gaussian negative
	log-likelihood \citep{KendallGal2017},
	\begin{equation}
		\ell(\theta,\phi) \;=\; \frac{1}{n}\sum_{i=1}^{n}
		\left[\frac{\big(\widehat{Y}_i - f_\theta(x_i)\big)^2}{2\,\sigma^2_\phi(x_i)}
		+ \tfrac12 \log \sigma^2_\phi(x_i)\right],
	\end{equation}
	so that the total predictive variance is
	$s^2(x) = s^2_{\mathrm{epi}}(x) + \sigma^2_\phi(x)$.

	\subsubsection{Finite-Sample Coverage}
	The point predictor and scale are wrapped in a split-conformal layer
	\citep{Lei2018}. On a held-out calibration set of size $m$, nonconformity scores
	$R_i = |\widehat{Y}_i - \bar f(x_i)|/s(x_i)$ are computed, and the interval at a
	new state is
	\begin{equation}
		C_\alpha(x) \;=\; \bar f(x) \;\pm\; q_{1-\alpha}\, s(x),
		\qquad
		q_{1-\alpha} = \big\lceil (m+1)(1-\alpha)\big\rceil\text{-th smallest } R_i .
	\end{equation}
	Under exchangeability of the scores the marginal coverage of the \emph{label}
	satisfies $\mathbb{P}\big(\widehat{Y}\in C_\alpha(x)\big)\ge 1-\alpha$. The
	target of inference, however, is the true value $Y(x)$, which the label
	estimates with variance $\sigma^2(x)/M$ from \eqref{eq:label}; certifying
	coverage of $Y(x)$ therefore requires accounting for that label noise, either by
	inflating $M$ so the noise is negligible relative to interval width or by
	deconvolving the known label variance from the calibration residuals. This
	label-noise correction is treated as part of the contribution, and is stated as
	a target result.
	\begin{proposition}[Target coverage, to be established]
		Under exchangeability of the sampled states and a stated control on the label
		noise $\sigma^2(x)/M$, the corrected split-conformal interval $C'_\alpha(x)$
		satisfies $\mathbb{P}\big(Y(x)\in C'_\alpha(x)\big)\ge 1-\alpha - o(1)$ as the
		inner-path count grows, with the deficit controlled explicitly by the label
		variance.
	\end{proposition}
	A key structural point is that exchangeability holds by construction here,
	because the training and calibration states are sampled independently from the
	design, in contrast to the dependent time-series setting where conformal
	coverage requires additional assumptions.

	\subsubsection{Experimental Design}
	States are drawn from a space-filling design over $\mathcal{X}$ so that accuracy
	and coverage can be assessed uniformly rather than only near a nominal contract.
	The inner-path count $M$, the training-set size $n$ and the network capacity are
	varied systematically to expose the label-budget to interval-width trade-off
	that RQ4 targets.

	\subsubsection{Validation and Capital Study}
	Values and intervals are benchmarked against full nested Monte Carlo and against
	the LSMC prices of \citet{LiLyu2026} and \citet{LangreneEtAl2026} on comparable
	contracts. Calibration quality is assessed by empirical coverage, mean interval
	width, the continuous ranked probability score and reliability diagrams.
	Reserves and a solvency-capital proxy are then computed with and without the
	quantified uncertainty to show its practical effect.

	\subsubsection{Empirical Application}
	Mortality is calibrated to the Human Mortality Database, with an optional South
	African case study substituting StatsSA or ASSA tables, which lets the thesis
	speak to a data-sparse and structurally distinct mortality environment.

	% ============================================================
	\section{Validity and Reliability} % Approx. 0.5 page

	\subsection{Validity}
	Accuracy is judged against nested Monte Carlo ground truth rather than against
	the emulator's own training labels, and coverage statements are made explicitly
	marginal or conditional so that what is claimed is what is tested. Counterexample
	contracts are retained to probe where the emulator degrades.

	\subsection{Reliability}
	Reproducibility is built in: versioned code, logged nested-simulation settings
	and random seeds, Monte Carlo error estimates on every reported number, and
	notebooks that regenerate each figure and table.

	% ============================================================
	\section{Ethical Considerations} % Approx. 0.5 page

	The research uses only publicly available, aggregated data (the Human Mortality
	Database and, where applicable, published South African mortality tables)
	together with synthetic contract portfolios. No individual policyholder records
	are used, so no participant consent arises. The applicable University of the
	Free State ethics or exemption process will be confirmed before data work
	begins, and data licences and attribution will be documented.

	% ============================================================
	\section{Expected Contributions of the Research} % Approx. 0.75 page

	\subsection{Theoretical Contributions}
	A finite-sample coverage guarantee for a deep surrogate of an actuarial
	valuation, including the treatment of Monte Carlo label noise and the
	exchangeability-by-construction argument, and a principled epistemic and
	aleatoric decomposition of the emulator's predictive uncertainty for this
	problem. A stretch result links the label budget, the network capacity and the
	resulting interval width. All results are labelled proposed until established.

	\subsection{Practical Contributions}
	A calibrated, uncertainty-aware VA emulator that values heterogeneous contracts
	in a single forward pass while reporting how far each value can be trusted,
	directly usable for pricing, hedging and solvency reporting, together with an
	open-source implementation and reproducible benchmarks against nested Monte
	Carlo and the 2026 LSMC methods.

	% ============================================================
	\section{Limitations of the Study} % Approx. 0.5 page

	The guarantees of split-conformal prediction are marginal rather than
	conditional, and hold under exchangeability of the sampled states; departures
	from the assumed data-generating model, the restriction to a chosen family of
	contracts, and the quality of the mortality calibration all bound what can be
	claimed. The emulator inherits any bias of the nested-simulation labels, and the
	South African application is constrained by data availability.

	% ============================================================
	\section{Timeline} % Approx. 0.5 page

	\begin{table}[H]
		\centering
		\caption{Indicative full-time schedule over three years.}
		\label{tab:timeline}
		\begin{tabular}{@{}ll@{}}
			\toprule
			\textbf{Months} & \textbf{Work} \\
			\midrule
			1 to 6    & Literature review, data audit, nested Monte Carlo label pipeline \\
			7 to 12   & Baseline emulator and accuracy study (C1); first paper \\
			13 to 20  & Uncertainty quantification and epistemic/aleatoric split (C2) \\
			21 to 24  & Conformal coverage theory and experiments (C3); second paper \\
			25 to 30  & Sample-complexity relations (C4); capital-impact study \\
			31 to 36  & South African case study, synthesis, third paper, examination \\
			\bottomrule
		\end{tabular}
	\end{table}

	Month 6 is the first decision point: confirm the label pipeline, the contract
	family and at least one working emulator before expanding the uncertainty
	analysis.

	% ============================================================
	\section{Budget} % Approx. 0.5 page

	The principal costs are GPU compute for nested-simulation label generation and
	emulator training, and access to market and mortality data. HMD data are free;
	GPU costs are modest on rented cloud instances, as in the candidate's MSc. A
	detailed budget will be priced before final submission.

	% ============================================================
	\section{Computational Resources} % Approx. 0.5 page

	Implementation uses Python and PyTorch with GPU acceleration, Bayesian
	hyperparameter optimisation, and a seeded, logged nested Monte Carlo label
	generator. A workstation with a modern GPU, supplemented by rented cloud GPUs
	for large label runs, is sufficient; no cluster access is assumed.

	% ============================================================
	\section{Chapter Layout} % Approx. 1 page

	The thesis is planned in seven chapters. Chapter 1 introduces the problem and
	context. Chapter 2 reviews VA guarantees, nested valuation, emulation,
	stochastic mortality and uncertainty quantification. Chapter 3 develops the
	emulator, its inputs and the nested Monte Carlo label generator. Chapter 4
	presents the uncertainty quantification and the epistemic and aleatoric
	decomposition. Chapter 5 establishes the finite-sample conformal coverage,
	including the label-noise correction. Chapter 6 reports the numerical validation
	and the capital-impact study, with the South African case. Chapter 7 draws
	conclusions, states limitations and sets out further work. The chapter count is
	an organisational choice, not a claim about depth.

	% ============================================================
	% References
	\clearpage
	\phantomsection
	\addcontentsline{toc}{section}{References}
	\nocite{*}
	\bibliographystyle{abbrvnat}
	\bibliography{../../References/References}

\end{document}
